\documentclass[10pt,a4paper]{article}
\usepackage[spanish,es-nodecimaldot]{babel}
\usepackage[utf8]{inputenc}
\usepackage[T1]{fontenc}
\usepackage{lmodern}
\usepackage[top=1.1cm,bottom=1.1cm,left=1.5cm,right=1.5cm]{geometry}
\usepackage{amsmath,amssymb}
\usepackage[table]{xcolor}
\usepackage{booktabs,array,multirow}
\usepackage{enumitem}
\usepackage{pgfplots}
\pgfplotsset{compat=1.18}
\usetikzlibrary{shapes.geometric,arrows.meta,positioning}

\setlength{\parindent}{0pt}
\setlength{\parskip}{2pt}
\pagestyle{empty}

\pgfmathdeclarefunction{gauss}{3}{%
  \pgfmathparse{1/(#3*sqrt(2*pi))*exp(-((#1-#2)^2)/(2*#3^2))}%
}

\begin{document}

{\large\bfseries TEMA A\par}
\vspace{2pt}

\textbf{Hipótesis nula ($H_0$).} Es el punto de partida y se asume temporalmente para confrontarla con los datos. Contiene la igualdad: $=$, $\leq$ o $\geq$.

\textbf{Hipótesis alternativa ($H_a$ o $H_1$).} Expresa lo que se concluirá si existe evidencia suficiente para rechazar $H_0$. Se formula con $<$, $>$ o $\neq$.

\vspace{3pt}
\textbf{Procedimiento prueba hipotesis}

\begin{enumerate}[label=\textbf{Paso \arabic*.},leftmargin=2cm,itemsep=0pt,topsep=1pt,parsep=0pt]
  \item Preguntar qué quiere demostrar el problema.
  \item Formular $H_0$ y $H_1$.
  \item Mirar el signo de $H_1$ y determinar si la prueba es derecha, izquierda o bilateral.
  \item Fijar el nivel de significancia $\alpha$.
  \item Elegir el estadístico: $Z$, $t$, $Z_{\bar{X}_1-\bar{X}_2}$, $t_{\bar{X}_1-\bar{X}_2}$ u otro que corresponda: $z_{\text{calc}} = \frac{\bar{x}-\mu_0}{\sigma/\sqrt{n}}$.
  \item Encontrar el valor crítico y construir la región de rechazo.
  \item Calcular $z_{calc}$ o $t_{calc}$ con los datos de la muestra.
  \item Comparar el estadístico calculado con el valor crítico, o comparar el $p$-valor con $\alpha$.
  \item Decidir: rechazar $H_0$ o no rechazar $H_0$.
  \item Interpretar la decisión con las palabras del problema.
\end{enumerate}

\vspace{3pt}
\textbf{Eleccion estadistico}

\begin{center}
\small
\renewcommand{\arraystretch}{1.2}
\begin{tabular}{|p{2.6cm}|p{3.1cm}|p{4.2cm}|p{4.2cm}|}
\hline
Población & Tamaño de la muestra & $\sigma$ conocida & $\sigma$ desconocida \\
\hline
\multirow{2}{=}{Normalmente Distribuida}
  & Grande ($n \geq 30$) & \centering $\bar{x} \pm z_\alpha \cdot \sigma_{\bar{x}}$ & \centering\arraybackslash $\bar{x} \pm z_\alpha \cdot s_{\bar{x}}$ \\
\cline{2-4}
  & Chica ($n < 30$)     & \centering $\bar{x} \pm z_\alpha \cdot \sigma_{\bar{x}}$ & \centering\arraybackslash $\bar{x} \pm t_\alpha \cdot s_{\bar{x}}$ \\
\hline
\multirow{2}{=}{No Normalmente Distribuida}
  & Grande ($n \geq 30$) & \centering $\bar{x} \pm z_\alpha \cdot \sigma_{\bar{x}}$ & \centering\arraybackslash $\bar{x} \pm z_\alpha \cdot s_{\bar{x}}$ \\
\cline{2-4}
  & Chica ($n < 30$)     & $\bar{x} \pm k \cdot \sigma_{\bar{x}}$ \newline donde $1-1/k^2$ se define por medio del Teorema de Chebyshev
                         & $\bar{x} \pm k \cdot s_{\bar{x}}$ \newline donde $1-1/k^2$ se define por medio del Teorema de Chebyshev \\
\hline
\end{tabular}
\end{center}

\vspace{-4pt}
\begin{center}
\small
\renewcommand{\arraystretch}{1.2}
\begin{tabular}{|c|c|c|}
\hline
 & Se Acepta H & Se Rechaza H \\
\hline
H es Verdadera & Decisión Correcta & Error de Tipo I \\
\hline
H es Falsa & Error de Tipo II & Decisión Correcta \\
\hline
\end{tabular}
\end{center}

\textbf{fabricante afirma q el secado es 20m, el comprador pinta tableros y lo rechaza si es mayor a 20.75min}

\begin{center}
\begin{tikzpicture}
\begin{axis}[
    width=12.5cm, height=6.2cm,
    xmin=18.8, xmax=23.2,
    ymin=0, ymax=1.0,
    axis x line=bottom,
    axis y line=left,
    xtick={19,20,21,22,23},
    ytick={0.2,0.4,0.6,0.8},
    clip=false,
    tick label style={font=\footnotesize},
    legend style={at={(0.99,0.32)},anchor=south east,draw=none,fill=none,font=\footnotesize},
    legend cell align=left,
]
% Area amarilla: error tipo II (bajo la curva azul, a la izquierda del limite)
\addplot[domain=19.3:20.75, samples=100, draw=none, fill=yellow!85!orange, forget plot] {gauss(x,21,0.5)} \closedcycle;
% Area gris: error tipo I (bajo la curva roja, a la derecha del limite)
\addplot[domain=20.75:22.2, samples=100, draw=none, fill=gray!55, forget plot] {gauss(x,20,0.5)} \closedcycle;
% Curvas
\addplot[red, thick, domain=18.8:23.2, samples=200] {gauss(x,20,0.5)};
\addlegendentry{$\mathrm{dnorm}\!\left(x,\mu,\frac{\sigma}{\sqrt{n}}\right)$}
\addplot[blue, thick, domain=18.8:23.2, samples=200] {gauss(x,21,0.5)};
\addlegendentry{$\mathrm{dnorm}\!\left(x,21,\frac{\sigma}{\sqrt{n}}\right)$}
% Limite de decision
\draw[very thick] (axis cs:20.75,0) -- (axis cs:20.75,0.92);
% Etiquetas de las medias
\node[above, font=\footnotesize] at (axis cs:20,0.82) {20};
\node[above, font=\footnotesize] at (axis cs:21,0.82) {21};
\node[above, font=\footnotesize] at (axis cs:20.75,0.93) {A};
% Error tipo II (beta)
\node[align=left, anchor=south west, font=\footnotesize] at (axis cs:18.85,1.10) {Probabilidad de Aceptar algo Incorrecto};
\node[draw, ellipse, font=\footnotesize, anchor=north west, inner sep=1.5pt] (t2) at (axis cs:18.85,1.08) {Error tipo II (beta)};
\draw[thick] (t2.south) -- (axis cs:20.55,0.12);
% Error tipo I (alfa)
\node[draw, ellipse, font=\footnotesize, anchor=west, inner sep=1.5pt] (t1) at (axis cs:21.6,1.02) {Error Tipo I (alfa)};
\node[align=left, anchor=north west, font=\footnotesize] at (axis cs:21.75,0.88) {Probabilidad de\\Rechazar algo\\correcto};
\draw[thick] (t1.south west) -- (axis cs:20.85,0.05);
\end{axis}
\end{tikzpicture}
\end{center}

\vspace{-4pt}
\begin{itemize}[leftmargin=1cm,itemsep=0pt,topsep=1pt,parsep=0pt]
  \item 20,75 es mi límite (elegido arbitrariamente)
  \item $\leftarrow$ hacia la izquierda del límite se acepta \\
        $\rightarrow$ hacia la derecha del límite se rechaza
  \item El area gris es la probabilidad de rechazar algo correcto (error tipo I)
  \item El area amarillo es la probabildiad de aceptar algo incorrecto (error tipo II)
\end{itemize}

\textbf{En el caso de q no lo elija arbitrariamente recurro al nivel de significancia $\alpha$}

\end{document}
